\documentclass[UTF8,a4paper,11pt,openany]{ctexbook}
\usepackage{../common/statlearnbook}
\SLSetBookMetadata{第 5 册}{统计学习方法讲义 v2.6.0}{采样方法、主题模型与图排序}
\begin{document}
\setcounter{page}{702}
\setcounter{chapter}{30}
\setcounter{figure}{2}
\pagestyle{headings}
\chaptermark{马尔可夫链基础}
\sectionmark{遍历定理与可逆性}
\paragraph{首次到达与首次正回返。}定义
\[
  \tau_i=\inf\{t\ge0:X_t=i\},\qquad
  \tau_i^+=\inf\{t\ge1:X_t=i\}.
\]
从状态 $i$ 出发，若 $\mathbb E_i[\tau_i^+]<\infty$，称 $i$ 正常返。图\ref{fig:V5-C01-return-time}强调正回返从 $t>0$ 开始计时，并与从其他状态出发的首次到达区分。
% v2.3.1 figure UID: FIG-P552-01
% v2.6.0 reverse redraw for FIG-P552-01: shared timeline grid and explicit readable hierarchy.
\tikzset{slfig-FIG-P552-01/.style={font=\fontsize{9.4pt}{11.2pt}\selectfont,every path/.append style={line cap=round,line join=round}}}
\begin{figure}[htbp]
\centering
\begin{tikzpicture}[slfig-FIG-P552-01,x=1.18cm,y=1cm,>=Stealth,
  event/.style={circle,draw=SLGray,fill=SLBlue!5,minimum size=6mm,line width=.68pt,
    font=\fontsize{9.2pt}{11.0pt}\selectfont},
  hit/.style={event,draw=SLGold,fill=SLGold!14,line width=.92pt},
  timeline/.style={-{Stealth[length=1.9mm]},draw=SLTextGray,line width=.78pt},
  interval/.style={-{Stealth[length=2.1mm]},draw=SLGold,line width=1pt},
  interval label/.style={font=\fontsize{9.2pt}{11.0pt}\selectfont,fill=white,inner sep=1.1pt},
  tick label/.style={font=\fontsize{8.6pt}{10.1pt}\selectfont},
  tick mark/.style={draw=SLTextGray,line width=.5pt},
  title/.style={font=\fontsize{9.4pt}{11.2pt}\selectfont\bfseries},
  note/.style={font=\fontsize{8.8pt}{10.4pt}\selectfont,fill=white,inner sep=1.1pt},
  conclusion/.style={draw=SLRuleGray,rounded corners=2pt,fill=white,align=center,
    minimum width=30mm,minimum height=12mm,line width=.68pt,inner sep=2.2pt,
    font=\fontsize{9.2pt}{11.0pt}\selectfont},
]
  \node[anchor=east,title,text=SLBlue] at (-.35,1.25) {首次到达 $\tau_i$};
  \draw[timeline] (0,1.25)--(7.45,1.25) node[right] {$t$};
  \foreach \t/\s in {0/j,1/k,2/k,3/i,4/k,5/i,6/k,7/k}{
    \draw[tick mark] (\t,1.33)--(\t,1.17);
    \node[below=2mm,tick label] at (\t,1.17) {$\t$};
    \node[event] at (\t,1.78) {$\s$};
  }
  \node[hit] at (3,1.78) {$i$};
  \draw[interval] (0,2.30)--(3,2.30)
    node[midway,above,interval label] {$j\ne i$ 出发，首次到达：$\tau_i=3$};

  \node[anchor=east,title,text=SLTeal] at (-.35,-1.10) {首次正回返 $\tau_i^+$};
  \draw[timeline] (0,-1.10)--(7.45,-1.10) node[right] {$t$};
  \foreach \t/\s in {0/i,1/k,2/k,3/i,4/k,5/i,6/k,7/k}{
    \draw[tick mark] (\t,-1.02)--(\t,-1.18);
    \node[below=2mm,tick label] at (\t,-1.18) {$\t$};
    \node[event] at (\t,-.57) {$\s$};
  }
  \node[hit] at (3,-.57) {$i$};
  \node[anchor=west,note,text=SLGold] at (.15,-.05) {$t>0$ 才计入};
  \draw[interval] (0,.10)--(3,.10)
    node[midway,above,interval label] {从 $i$ 出发，首次正回返：$\tau_i^+=3$};

  \node[conclusion] at (8.85,.10)
    {正常返考察\\$\mathbb E_i[\tau_i^+]<\infty$};
\end{tikzpicture}
\caption{首次到达与首次正回返的区分：正常返考察从$i$出发的$\mathbb E_i[\tau_i^+]$，而不是某个固定时刻首次到达的概率}\label{fig:V5-C01-return-time}
\end{figure}

\end{document}
